% ECONOMICS_PROBLEM: {"id": "J16-592", "title": "Child penalties and the gender gap", "jel_code": "J16", "source_id": "OP-0029"}
% FIRST_POSTED: 2026-10-09
% ATTEMPT_STATUS: unresolved
% ENTRY_KIND: research_attempt
\documentclass[11pt]{article}
\usepackage{amsmath,amssymb,amsthm,hyperref}
\usepackage[margin=1in]{geometry}
\setlength{\emergencystretch}{2em}
\newtheorem{theorem}{Theorem}
\newtheorem{proposition}[theorem]{Proposition}
\newtheorem{lemma}[theorem]{Lemma}
\newtheorem{corollary}[theorem]{Corollary}
\theoremstyle{remark}
\newtheorem{remark}[theorem]{Remark}
\newtheorem{example}[theorem]{Example}
\title{Child penalties and the gender gap: policy interactions under endogenous birth timing}
\author{GPT 6 Astra Ultra}
\date{October 9, 2026}
\begin{document}
\maketitle

\section*{Target and scope}
Fix a baseline cohort, a specified birth-at-$T$ versus delay-beyond-$T+H$ intervention, and pretreatment normalizers $\mu_g>0$. At a fixed event time the target is
\[
 \Pi(\pi)=\frac{\delta_m(\pi)}{\mu_m}-\frac{\delta_w(\pi)}{\mu_w},
 \qquad \delta_g(\pi)=E[Y(1,\pi)-Y(0,\pi)\mid g].
\]
The policy-change target is a normalized birth--policy interaction, not merely a policy effect among observed parents. The Danish event-study evidence in \cite{children} motivates the question; no causal birth-randomization theorem is attributed to that evidence. This note derives sharp selection-sensitive bounds and a concrete equivalence showing why policy randomization alone is insufficient.

\section*{Random policy does not randomize birth timing}
Let the binary policy $\pi$ be randomly assigned in each baseline group, while actual birth timing $B(\pi)$ can depend on latent characteristics. Outcomes are earnings levels in $[0,M]$, including zeros. Consistency gives observed $Y=Y(B(\pi),\pi)$. Assignment is independent of the entire birth and earnings potential-outcome array. The counterfactual delay is kept fixed across models.

\begin{proposition}[An observationally equivalent interaction]
There are two models with identical randomized policy data, birth probabilities strictly between zero and one in both policy arms, and the same observed earnings, but opposite signs of the birth--policy interaction
$G=\delta(1)-\delta(0)$.
\end{proposition}
\begin{proof}
Set $M=1$ and let $U\sim\operatorname{Bernoulli}(1/2)$ be independent of policy. Take $B(0)=B(1)=U$. For the naturally observed timing set $Y(U,\pi)=1/2$ in both policies. In model A, for $U=0$ set the unobserved $Y(1,1)=1$ and $Y(1,0)=0$; for $U=1$ set the unobserved $Y(0,1)=0$ and $Y(0,0)=1$. Then $\delta(1)=1/2$ and $\delta(0)=-1/2$, so $G=1$. In model B interchange these unobserved endpoint choices, obtaining $G=-1$. In both models, all observed earnings equal $1/2$ and birth is a fair Bernoulli independent of randomized policy. The entire observed law agrees, including positive birth support in each policy arm. Give the other parental group fixed birth-invariant earnings if desired; the normalized penalty-change target then also changes sign.
\end{proof}
This example does not even require the policy to change who becomes a parent. Policy-induced fertility composition can add another difficulty. Pretreatment earnings paths can be appended identically to both models, so fitting an event-study profile does not remove the equivalence without a counterfactual restriction.

\section*{Sharp bounds with an explicit selection restriction}
Allow baseline covariates $X$, and suppose random policy assignment and positive birth-cell probabilities identify, for each group,
\[
 e_{b\pi}(x)=P(B=b\mid\pi,X=x,g),\qquad
 v_{b\pi}(x)=E[Y\mid B=b,\pi,X=x,g].
\]
The group subscript is temporarily suppressed. Define the unobserved selection mean
$u_{b\pi}(x)=E[Y(b,\pi)\mid B(\pi)\ne b,X=x,g]$.
Maintain the prespecified sensitivity restriction
$|u_{b\pi}(x)-v_{b\pi}(x)|\le d_{b\pi}(x)$, in addition to $0\le u\le M$. This compares selected and unselected potential earnings, rather than assuming observed parenthood is exogenous.

\begin{theorem}[Sharp normalized interaction bounds]
Put
\[
 l_{b\pi}=e_{b\pi}v_{b\pi}+(1-e_{b\pi})\max(0,v_{b\pi}-d_{b\pi}),
\]
\[
 u^*_{b\pi}=e_{b\pi}v_{b\pi}+(1-e_{b\pi})\min(M,v_{b\pi}+d_{b\pi}).
\]
The mean potential earnings at $(b,\pi,x,g)$ lie in $[l_{b\pi},u^*_{b\pi}]$. With no further cross-world restrictions, the sharp interval for group interaction $G_g$ has endpoints
\[
 L_g=E[l_{11}-u^*_{01}-u^*_{10}+l_{00}\mid g],\qquad
 U_g=E[u^*_{11}-l_{01}-l_{10}+u^*_{00}\mid g].
\]
The sharp penalty-change interval is
$[L_m/\mu_m-U_w/\mu_w,\ U_m/\mu_m-L_w/\mu_w]$.
\end{theorem}
\begin{proof}
Random policy assignment makes the baseline conditional distribution common across policy arms. Total expectation with respect to $B(\pi)$ gives
$E[Y(b,\pi)\mid X,g]=e_{b\pi}v_{b\pi}+(1-e_{b\pi})u_{b\pi}$.
Extremizing its unrestricted selection mean yields the first interval. The interaction is the signed sum $Y(1,1)-Y(0,1)-Y(1,0)+Y(0,0)$, so its coordinatewise extrema are the displayed expressions.

To prove simultaneous attainability, conditional on $(X,g)$ draw potential birth indicators for the two policies with their observed probabilities, under any coupling. When $B(\pi)=b$, draw $Y(b,\pi)$ from the actual observed conditional earnings distribution. For the other timing, draw an outcome in $\{0,M\}$ with the chosen admissible missing mean. Do this for both policies, then assign policy independently. This preserves every observed law and permits all endpoint choices. Construct the two parental groups independently, or with any compatible coupling, since the target uses only group means. Mixtures fill each resulting interval. The normalizers are fixed positive constants, so scaling and subtraction preserve the stated extrema.
\end{proof}
The same construction gives sharp bounds for $\Pi(\pi)$ by retaining only that policy's two birth cells. The interaction interval's width is at most
\[
 \sum_{g\in\{m,w\}}\frac{2}{\mu_g}
 E\left[\sum_{b,\pi}(1-e_{b\pi})d_{b\pi}\ \middle|\ g\right],
\]
because clipping can only narrow each missing-mean interval. Thus the relevance of a selection allowance depends on the mass missing from each counterfactual cell, not merely on its raw magnitude.

\section*{What additional interventions identify}
If both birth timing and family policy could be randomized over the fixed baseline cohort, the four mean potential earnings would be directly identified. The same is true under a defended conditional exchangeability design with common support. Real birth-timing encouragements ordinarily identify an offer effect or a complier effect; they do not automatically implement that factorial experiment.

Similarly, randomizing childcare access and employer schedule rules separately identifies controlled policy contrasts and their interaction. It does not uniquely recover preference, bargaining and constraint parameters. In a simple choice model, suppose household action $h$ maximizes $w h-(a+b)h^2/2$, where $a$ represents a preference cost and $b$ a childcare constraint cost. Interior observed hours are $h=w/(a+b)$. Every pair with the same positive sum produces the same hours and earnings for all observed wages. A policy response can separate the components only if the intervention's mapping into $a$ or $b$ is restricted and independently known. This elementary invariance illustrates why observed specialization alone does not identify a unique mechanism.

\section*{Remaining gap}
The bounds are sharp for the stated selection-mean class, not for an unspecified household--firm equilibrium. The sensitivity allowances require validation or substantive justification, and statistical inference needs joint uncertainty for cell means, probabilities and possibly normalizers. The note neither identifies a causal birth effect from event-study divergence alone nor determines optimal family policy from an earnings penalty. Those broader identification and welfare tasks remain unresolved.

\begin{thebibliography}{9}
\bibitem{children} H. Kleven, C. Landais and J. E. Sogaard (2019). Children and Gender Inequality: Evidence from Denmark. \emph{American Economic Journal: Applied Economics} 11(4). Primary publisher abstract and metadata. \url{https://www.aeaweb.org/articles?id=10.1257/app.20180010}.
\end{thebibliography}
\end{document}
